%=======================================================================
%  CSE558 Data Science - Midterm Project Report
%  Monsoon 2026, IIIT-Delhi
%  Built on CSE558_Project_Report_Template.tex. Numbers come from
%  ../*.md and ../*.json, produced by the scripts in ../../src/.
%  Compile with pdfLaTeX, twice, from this folder.
%=======================================================================
\documentclass[10pt,a4paper,twocolumn]{article}

%========================= EDIT THIS =========================
\newcommand{\projecttitle}{Same job, higher wage? E-3 versus H-1B filings in
  US Labor Condition Applications, FY2017--2022}
\newif\ifshowhints
\showhintsfalse
%=============================================================

%-------------------------------------------- fonts (no need to edit)
\usepackage{iftex}
\ifPDFTeX
  \usepackage[T1]{fontenc}
  \usepackage[utf8]{inputenc}
  \IfFileExists{carlito.sty}%
    {\usepackage[sfdefault]{carlito}}%              Calibri look
    {\usepackage[scaled=0.92]{helvet}\renewcommand{\familydefault}{\sfdefault}}%
  \usepackage{courier}
\else
  \usepackage[no-math]{fontspec}
  \setmainfont{Carlito}[Extension=.ttf, UprightFont=*-Regular, BoldFont=*-Bold,
    ItalicFont=*-Italic, BoldItalicFont=*-BoldItalic]
\fi

%----------------------------------------------------------- packages
\usepackage[margin=18mm,bottom=20mm,columnsep=7mm]{geometry}
\usepackage{amsmath}
\usepackage{xcolor}
\usepackage{graphicx}
\usepackage{tabularx}
\usepackage{array}
\usepackage{booktabs}
\usepackage{enumitem}
\usepackage{titlesec}
\usepackage{needspace}
\usepackage[font=small,labelfont=bf]{caption}
\usepackage[hidelinks]{hyperref}
\hypersetup{pdftitle={\projecttitle},
            pdfsubject={CSE558 Data Science, Monsoon 2026: Midterm Project Report}}

%------------------------------------------------------------- layout
\definecolor{eg}{HTML}{7396C4}
\definecolor{hintgrey}{gray}{0.58}
\definecolor{todo}{HTML}{C0392B}       % placeholders still to fill in

\titleformat{\section}{\normalfont\large\bfseries}{\thesection.}{0.5em}{}
\titlespacing*{\section}{0pt}{12pt}{3pt}
\titleformat{\subsection}{\normalfont\normalsize\bfseries}{\thesubsection}{0.5em}{}
\titlespacing*{\subsection}{0pt}{10pt}{2pt}
\newcommand{\sectionbreak}{\needspace{10\baselineskip}}
\newcommand{\subsectionbreak}{\needspace{10\baselineskip}}

\setlength{\parindent}{0pt}
\setlength{\parskip}{4pt}
\setlist{leftmargin=1.2em, itemsep=2pt, topsep=2pt, parsep=0pt}
\renewcommand{\arraystretch}{1.25}
\setlength{\tabcolsep}{4pt}
\newcolumntype{L}[1]{>{\raggedright\arraybackslash}p{#1}}
\newcolumntype{Y}{>{\raggedright\arraybackslash}X}
\newcolumntype{K}[1]{>{\raggedright\arraybackslash\bfseries\boldmath}p{#1}}

%--------------------------------------------------------- shorthands
\newcommand{\hint}[1]{\ifshowhints{\footnotesize\itshape\color{hintgrey}#1\par}\fi}
\newcommand{\field}[2]{\textbf{#1}~#2\par}
\newcommand{\tpart}[2]{\par\needspace{12\baselineskip}\smallskip\textbf{#1)~#2}\par}
\newcommand{\minihead}[1]{\par\needspace{8\baselineskip}\smallskip\textbf{#1}\par}
\newcommand{\code}[1]{\texttt{#1}}
\newcommand{\todo}[1]{\textcolor{todo}{[#1]}}
\newcommand{\figfile}[2]{%
  \IfFileExists{#1}%
    {\includegraphics[width=\linewidth,height=#2,keepaspectratio]{#1}}%
    {\fbox{\parbox[c][#2][c]{\dimexpr\linewidth-2\fboxsep-2\fboxrule\relax}%
      {\centering\small\itshape\color{hintgrey}Your plot: \detokenize{#1}}}}}
\newcommand{\figdir}{../../figures/}

%=======================================================================
\begin{document}

\twocolumn[%
{\small\bfseries CSE558 Data Science\ \ \textbullet\ \ Midterm Project Report\ \ \textbullet\ \ Monsoon 2026\par}
\vspace{4pt}
{\LARGE\bfseries\raggedright\projecttitle\par}
\vspace{6pt}
\rule{\linewidth}{1pt}

{\small
\begin{tabularx}{\linewidth}{@{}K{2.4cm} Y@{}}
\toprule
Group name & Ramanujan et al. \\
Dataset & kaggle.com/datasets/jishnukoliyadan/lca-programs-h1b-h1b1-e3-visa-petitions \\
Code & \url{https://github.com/sjwl07/h1b-e3-wage-gap} \\
\bottomrule
\end{tabularx}}

\vspace{2pt}
{\small
\begin{tabularx}{\linewidth}{@{}L{0.5cm} L{3.4cm} L{2.3cm} L{4.8cm} Y@{}}
\toprule
\textbf{\#} & \textbf{Name} & \textbf{Roll no.} & \textbf{Institute email} & \textbf{Contribution} \\
\midrule
1 & Shivam Jhunjhunuwala & 2023505 & shivam23505@iiitd.ac.in & Data loading, held-out split, report \\
2 & Ramanujan Gopalan & 2023424 & ramanujan23424@iiitd.ac.in & Data dictionary, cleaning rules \\
3 & Saanvi Singh & 2023454 & saanvi23454@iiitd.ac.in & EDA figures and captions \\
4 & Mohammad Umar & 2023324 & umar23324@iiitd.ac.in & Research questions and test planning \\
5 & Mayank Kumar & 2023317 & mayank23317@iiitd.ac.in & Hypothesis tests, robustness checks \\
\bottomrule
\end{tabularx}}
\vspace{12pt}
]

%-----------------------------------------------------------------------
\section{Dataset}

\field{Source:}{Kaggle, ``LCA Programs (H-1B, H-1B1, E-3) Visa Petitions'' (CC BY-NC-SA 4.0), an
extract of the US Department of Labor's Labor Condition Application (LCA) disclosure data;
six yearly CSV files, FY2017--2022, downloaded on 21 October 2026. One row is one LCA: an
employer's filing, required before hiring a foreign worker on an H-1B, H-1B1 or E-3 visa, that
states the job, the worksite and the \emph{prevailing wage} (the legal minimum DOL sets for that
occupation, area and wage level). DOL then certifies or denies it.}

\field{What is raw about it:}{the case status is spelled differently before and after FY2020;
FY2019's wage column holds the \emph{offered} wage, not the prevailing wage; FY2019's worksite is
blank in 88\% of rows; wages come in five pay units with typing errors (an ``hourly'' wage of
46{,}301, an annual wage of \$1\,billion); 1{,}769 spellings for about 1{,}100 occupation (SOC)
titles; 255{,}125 spellings of employer names; 24\% of rows are exact copies of another row.}

{\small
\begin{tabularx}{\linewidth}{@{}L{2.9cm} Y Y@{}}
\toprule
 & \textbf{Raw files} & \textbf{After encoding} \\
\midrule
Instances (rows) & 3{,}973{,}349 & 3{,}178{,}679 \\
Attributes (columns) & 11--12 & 210 \\
Instances $\times$ attributes & 47{,}680{,}188 & 667{,}522{,}590 \\
\bottomrule
\end{tabularx}}

\field{Size check:}{the encoded figure is for the 80\% exploration set alone (the 20\%
confirmation set is held out, Section~2): $3{,}178{,}679\times210=667{,}522{,}590\ge3{,}000{,}000$.
The 210 features are 197 one-hot levels (the 50 most common SOC titles plus \code{Other},
worksite and employer state, year, quarter, visa class, status, pay unit, full-time, region)
and 13 numeric features and flags. The raw files alone already pass:
$3{,}973{,}349\times12\approx47.7$M.}

\minihead{Key features} (every attribute's type, unit, range and missingness by year: \code{report/data\_dictionary.md})
{\footnotesize
\begin{tabularx}{\linewidth}{@{}L{2.65cm} L{1.5cm} Y@{}}
\toprule
\textbf{Feature} & \textbf{Type} & \textbf{Real-world meaning} \\
\midrule
\code{log\_pw} & continuous & log of the annual prevailing wage (\$15k--\$462k after cleaning); response for P1 \\
\code{case\_status} & nominal & DOL decision: certified (90.9\%), certified-withdrawn, withdrawn, denied (0.9\%); \code{denied} is the response for P2 \\
\code{visa\_class} & nominal & H-1B (97.6\%), E-3 Australian (2.0\%), H-1B1 Chile/Singapore (0.4\%) \\
\code{soc\_title} & nominal & occupation, 1{,}123 titles after cleaning; 30\% are ``Software Developers, Applications'' \\
\code{worksite\_state} & nominal & where the worker will work, 55 states and territories \\
\code{fiscal\_year} & discrete & US fiscal year of the decision, 2017--2022 \\
\code{employer\_key} & nominal (ID) & normalised employer name: 165{,}799 in all six years, 132{,}485 in the exploration set; used for grouping, not encoded \\
\code{employer\_volume} & discrete & derived: filings by that employer that year, 1--21{,}448 \\
\code{full\_time} & nominal & full-time position (98.2\%) \\
\bottomrule
\end{tabularx}}

%-----------------------------------------------------------------------
\section{Data preparation}
All rules were fitted on the exploration set and saved; the confirmation set was later cleaned
with the same saved rules, never re-fitted. No rows are dropped: every row keeps a valid case
decision, so a bad wage is set to missing and flagged instead.

{\footnotesize
\begin{tabularx}{\linewidth}{@{}L{1.95cm} L{2.35cm} Y@{}}
\toprule
\textbf{Issue found} & \textbf{What we did} & \textbf{Why, and not the alternative} \\
\midrule
hypotheses would come from EDA
 & split 80/20 \emph{before} any EDA: stratified group k-fold, grouped by employer, stratified by year $\times$ status, seed 42
 & testing on the data that suggested a hypothesis is circular. Grouping keeps one employer's near-identical filings out of both sets; a random split would leak them \\
status spelled two ways (upper case before FY2020, title case after)
 & mapped to 4 labels
 & otherwise the same decision counts as two levels \\
\bottomrule
\end{tabularx}}

{\footnotesize
\begin{tabularx}{\linewidth}{@{}L{1.95cm} L{2.35cm} Y@{}}
\toprule
\textbf{Issue found} & \textbf{What we did} & \textbf{Why, and not the alternative} \\
\midrule
FY2019 wage column is the offered wage
 & kept as a separate column; FY2019's prevailing wage left missing
 & 39\% of FY2019 values are whole thousands, against under 1\% in every other year, and the median per SOC title jumps $+14.5\%$ then falls $-8\%$. Merging would shift FY2019 wages up about 10\% \\
777{,}688 exact duplicate rows (24\%)
 & \textbf{kept}, counted in \code{n\_identical}
 & there's no case number, and in DOL's source every row is a separate LCA. The duplicate share rises with employer size (18\% $\to$ 55\%) and differs by status (36\% certified vs 16\% withdrawn), which an export glitch wouldn't cause. Dropping would delete real filings of big employers; their correlation is handled by clustering (Section~5) \\
wrong pay unit (``Hour'' $=46{,}301$)
 & 427 relabelled to Year
 & rule: value $\times$ unit $>\$1$M/year but the value itself is a normal salary (\$15k--\$1M); fixing keeps a valid wage \\
five pay units
 & annualised ($\times$2080, 52, 26, 12)
 & puts all wages on one scale (2{,}080\,h = DOL's full-time year) \\
\bottomrule
\end{tabularx}}

{\footnotesize
\begin{tabularx}{\linewidth}{@{}L{1.95cm} L{2.35cm} Y@{}}
\toprule
\textbf{Issue found} & \textbf{What we did} & \textbf{Why, and not the alternative} \\
\midrule
101 wages of \$0 or with no unit
 & set to missing
 & impossible values; the row's decision is still valid \\
1{,}769 SOC title spellings (codes glued on, typos, plurals)
 & merged to 1{,}123 titles: letters-only key without plurals, then fuzzy match (similarity $\ge0.92$) of rare spellings
 & exact matching leaves e.g.\ ``15-1121.00Computer\-systems\-analysts'' as a separate job. We checked the merges by hand; a looser cutoff wrongly joined ``Electrical'' and ``Electronics'' engineers \\
wage outliers (\$1B ``annual'' wage)
 & IQR fence (1.5) on \emph{log} wage \emph{within} SOC title; 21{,}579 wages set to missing and flagged
 & log, because wages are right-skewed; within title, because \$200k is normal for a surgeon, not a cook. Setting the wage (not the row) to missing keeps the case decision \\
employer country blank (100{,}916) or ``Georgia''
 & set to USA when the employer state is a US state
 & the address shows the country; a US state was typed as the country \\
FY2019 worksite 88\% blank
 & flagged, not imputed
 & worksite and employer state agree in only 52\% of rows, and the blank rows differ (93\% vs 61\% certified), so imputing would invent locations \\
columns to encode
 & year, quarter, status, visa, pay unit, full-time, states, region one-hot; flags 0/1; employer volume logged
 & one-hot assumes no order, so wages may move non-linearly by year; the two responses stay numeric (\code{log\_pw}, \code{denied} 0/1) for the tests; volume is very skewed (1--21{,}448) \\
high-cardinality text
 & one-hot only 10 low-cardinality columns; SOC as top 50 + \code{Other} (87.7\% of rows); employer and job title not one-hot
 & 132k employers and 375k job titles would add half a million near-empty columns; they enter as filing volume and seniority flags instead \\
\bottomrule
\end{tabularx}}

%-----------------------------------------------------------------------
\section{What changed after preparation}
Exploration set, before and after cleaning (\code{report/cleaning\_log.md}).

{\footnotesize
\begin{tabularx}{\linewidth}{@{}L{3.0cm} L{1.9cm} Y@{}}
\toprule
\textbf{Measure} & \textbf{Before} & \textbf{After} \\
\midrule
Rows & 3{,}178{,}679 & 3{,}178{,}679 (none dropped) \\
Attributes & 17 harmonised$^{*}$ & 40 cleaned; 210 after encoding \\
Usable prevailing wages & 2{,}645{,}476 & 2{,}623{,}796 (21{,}680 set missing: 21{,}579 outliers, 101 invalid) \\
Wage range (annual) & \$0 -- \$1B & \$15{,}080 -- \$462{,}342 \\
Pay units & 5, 427 mislabelled & 1 (annual) \\
SOC title spellings & 1{,}769 & 1{,}123 \\
Job title spellings & 439{,}495 & 375{,}194 \\
Employer name spellings & 204{,}209 & 132{,}485 employers \\
Blank employer country & 100{,}916 & 0 \\
Exact duplicate rows & 777{,}688 & 777{,}688, kept and counted \\
Missing cells (not by design) & 1.92\% & 1.56\% \\
\bottomrule
\end{tabularx}}
{\footnotesize $^{*}$the 12 raw columns, with worksite and employer location split into city and state, plus
year, employer key and FY2019's offered wage.}

Cleaning turned the wage into one comparable annual number: the lowest cleaned wage is exactly
$\$7.25\times2{,}080$, the US federal minimum wage, a sign the unit fixes worked. Its log is almost
symmetric (skew $-0.04$, against $1.14$ for dollars; Fig.~\ref{fig:clean}). Harmonising occupation and employer names
is what makes the within-occupation and within-employer comparisons of Section~5 possible.
\textbf{What we lost:} 21{,}680 wages (0.8\%) set to missing; all of FY2019 (17\% of rows) for
wage questions, as it has no prevailing wage; the detail of the 1{,}073 rarer SOC titles, grouped
as \code{Other} (12.3\% of rows) in the encoding; and any distinction between fuzzy-merged
spellings, if a merge was wrong.

\begin{figure*}[!t]
  \figfile{\figdir fig01_wage_distribution.png}{4.2cm}
  \caption{Before cleaning, wages span \$25 to \$1B a year (grey); after, \$15k--\$462k with a symmetric log (right).}
  \label{fig:clean}
\end{figure*}

%-----------------------------------------------------------------------
\section{Basic questions and problem statement}

\minihead{Basic questions}
\begin{enumerate}[label=\arabic*.]
  \item Which employers file the most? (Infosys, TCS, Google, Amazon, EY; the top 10 file 16\%; Fig.~\ref{fig:basic}.)
  \item Which jobs? (Mostly software: computer titles are 65\% of filings; Fig.~\ref{fig:soc}.)
  \item How many filings are certified? (91\%; only 0.87\% are denied.)
  \item Have wages risen? (Median prevailing wage $+33\%$ from FY2017 to FY2022; Fig.~\ref{fig:basic}.)
  \item Do some regions pay more? (West \$105k vs Midwest \$78k; Fig.~\ref{fig:basic}.)
\end{enumerate}

\begin{figure*}[!t]
  \figfile{\figdir fig15_basic_questions.png}{4.2cm}
  \caption{Basic questions 1, 4, 5: the top 10 employers file 16\%, the median wage rose 33\%, the West pays most.}
  \label{fig:basic}
\end{figure*}

\begin{figure*}[!t]
  \figfile{\figdir fig04_soc.png}{4.2cm}
  \caption{Software jobs dominate (one title is 30\% of filings) and wages vary within each title,
  so occupation must be held fixed.}
  \label{fig:soc}
\end{figure*}

\begin{figure*}[!t]
  \figfile{\figdir fig05_visa.png}{4.6cm}
  \caption{E-3 filings sit $+8.8\%$ above the median of the same SOC $\times$ state $\times$ year,
  H-1B $+1.6\%$, and the gap holds at every employer size.}
  \label{fig:visa}
\end{figure*}

\begin{figure*}[!t]
  \figfile{\figdir fig07_employer_volume.png}{4.6cm}
  \caption{Denials peak at employers with 2--10 filings a year, yet E-3 is denied more within every size group.}
  \label{fig:volume}
\end{figure*}

\begin{figure}[!ht]
  \figfile{\figdir fig09_seniority.png}{3.9cm}
  \caption{Title keywords move relative wage as expected (senior $+10\%$, junior $-16\%$), so it
  tracks the wage level an employer chose.}
  \label{fig:senior}
\end{figure}

\minihead{Our problem statements}
\noindent\fbox{\begin{minipage}{\dimexpr\linewidth-2\fboxsep-2\fboxrule\relax}
\textbf{P1 (primary).} For the same occupation, worksite state and fiscal year, do E-3
(Australian) filings carry a different prevailing wage than H-1B filings, and does any
difference remain within the same employer?

\smallskip
\textbf{P2.} Are E-3 filings denied more or less often than H-1B filings from employers of the
same size in the same year, and within the same employer?
\end{minipage}}

\field{How they are novel:}{the basic questions summarise one variable at a time. The prevailing
wage is set by DOL from occupation, area and wage level; the visa isn't an input, so ``no
difference'' is a natural default for P1. A difference would mean employers place the two groups
at different wage levels, which the data records only indirectly (Fig.~\ref{fig:senior}). P1 holds
fixed what the basic questions show moves wages most (occupation, place, year) and then the
employer. P2 asks whether the raw denial gap (E-3 denied 5.8$\times$ as often in the exploration
set, 6.9$\times$ in the confirmation set) survives the main confounder found in EDA, employer size
(Fig.~\ref{fig:volume}). EDA rules out a Simpson's paradox: E-3 is denied more within every size
group too, so the question is how much of the gap size explains.}

\field{Correlation, not causation:}{both statements are worded as differences between
\emph{filings}, not effects of the visa. Australians may simply be hired into more senior roles,
and worker experience and education aren't in the data. A causal claim would need worker-level
data to compare similar workers, or a rule change that affects only one visa, studied before and
after. Neither is available, so we make no causal claim.}

%-----------------------------------------------------------------------
\section{Statistical inferences}

\field{Rows tested:}{the held-out confirmation set: 794{,}670 filings from 33{,}314 employers not in
the exploration set, cleaned with the saved rules; E-3 vs H-1B only (H-1B1 is 0.4\%, a separate
program). Wage tests use FY2017--18 and FY2020--22 (FY2019 has no prevailing wage) and rows with a
valid wage and worksite.}

\field{Common design, fixed before the held-out data was opened:}{each test is $\beta$ in
$y_i=\beta\,\mathrm{E3}_i+\alpha_{c(i)}+\varepsilon_i$, where $\alpha_c$ is a separate intercept
(\emph{fixed effect}) for each comparison cell, e.g.\ SOC $\times$ state $\times$ year. Each choice
has one reason:}
\begin{itemize}[leftmargin=*,itemsep=1pt,topsep=1pt]
  \item \textbf{Cell intercepts}: $\beta$ becomes the average E-3 minus H-1B gap \emph{inside}
  cells, so differences in job, place and year can't create a gap.
  \item \textbf{Standard errors clustered by employer}: one employer's filings are near-copies
  (Section~2), so they aren't independent; treating them as independent makes the SE too small.
  The independent units are employers, so the $t$-test uses employers $-\,1$ degrees of freedom.
  \item \textbf{Holm correction}: four tests at 5\% each would give up to a 19\% chance of at least one
  false alarm; Holm keeps it at 5\% overall.
  \item \textbf{Smallest meaningful difference (MID)}: with this many rows even a trivial gap is
  significant, so each result is also judged by its 95\% CI: \emph{meaningful} (whole CI beyond
  the MID), \emph{possibly too small} (CI includes the MID) or \emph{negligible} (CI inside it).
\end{itemize}
Before the run we checked power, the chance of detecting a gap exactly the size of the MID:
T1 1.00, T2 0.999, T3 0.16, T4 0.93.

%------------------------------------------------------------ Test 1
% keeps page 4 for figures only, so Test 1 starts at the top of page 5; remove if text above changes
\mbox{}\newpage\mbox{}\newpage\mbox{}\newpage
\subsection*{Test 1 (P1): do E-3 filings carry a different prevailing wage for the same job,
state and year?}

\tpart{a}{Hypotheses, written before running the test}
{\small
\begin{tabularx}{\linewidth}{@{}K{3.4cm} Y@{}}
\toprule
$H_0$ & $\beta=0$ \\
$H_1$ & $\beta\neq0$ \\
Significance level & $\alpha=0.05$ family-wise (Holm, 4 tests) \\
Side & two-sided \\
Test & $t$-test on $\beta$, fixed effects for SOC $\times$ state $\times$ year, employer-clustered SE \\
Smallest difference that matters & 5\% ($|\beta|\ge\log1.05=0.049$) \\
\bottomrule
\end{tabularx}}

$y$ = log annual prevailing wage, so $\beta$ is the E-3 minus H-1B difference in log points
($\approx$\%). Two-sided, because P1 and P2 ask whether the groups differ in either direction. 5\%,
because it is about \$4{,}400 at the median wage, about one year of wage growth here, and half
the gap between ``senior'' and other titles in the same job: a smaller gap wouldn't place a
worker at a different wage level.

\tpart{b}{Other tests we considered}
Two groups and a continuous, symmetric response suggest a two-sample $t$-test, but the groups
must be compared within the same occupation, state and year (Figs.~\ref{fig:soc}, \ref{fig:visa}),
and filings aren't independent. A fixed-effects regression does both comparisons at once
for 2{,}729 cells, and clustering by employer handles the dependence.
{\footnotesize
\begin{tabularx}{\linewidth}{@{}L{2.4cm} Y@{}}
\toprule
\textbf{Also considered} & \textbf{Why we did not use it} \\
\midrule
Welch $t$-test on relative wage & treats filings as independent; its SE (0.0027) is 2.4$\times$ too small. Run as a check: $t=20.7$, same decision \\
Mann--Whitney U & tests a shift in ranks, which can't be compared with a 5\% MID; log wage is already symmetric, so ranks add nothing; also assumes independence \\
Raw two-group comparison & mixes job, state and year composition into the gap \\
\bottomrule
\end{tabularx}}

\tpart{c}{Result}
{\small
\begin{tabularx}{\linewidth}{@{}K{3.4cm} Y@{}}
\toprule
Rows used & 408{,}671 (10{,}433 E-3), 19{,}299 employers \\
Estimate & $+9.1\%$ ($\hat\beta=0.087$; raw gap $+9.6\%$) \\
95\% confidence interval & $+7.7\%$ to $+10.5\%$ \\
Test statistic (d.f.) & $t=13.66$ (19{,}298) \\
p-value (raw, Holm-adjusted) & $3.0\times10^{-42}$ ($8.9\times10^{-42}$) \\
Decision & reject $H_0$ \\
Verdict & meaningful: the whole CI is above 5\% \\
\bottomrule
\end{tabularx}}

For the same occupation, state and year, E-3 filings carry a prevailing wage about 9\% higher than
H-1B filings, well past the 5\% that matters. Occupation, state and year explain almost none of
the raw $+9.6\%$ gap. The estimate is stable across years ($+8$ to $+12\%$, Fig.~\ref{fig:year})
and unchanged without the largest employer.

\begin{figure*}[!t]
  \figfile{\figdir fig13_inference_estimates.png}{4.3cm}

  \figfile{\figdir fig14_inference_by_year.png}{4.3cm}
  \caption{Top: estimates with 95\% CIs (grey $=\pm$MID). Bottom: Tests 1, 3 by year (exploratory); the denial gap shrinks from FY2018.}
  \label{fig:year}
\end{figure*}

%------------------------------------------------------------ Test 2
\needspace{18\baselineskip}
\subsection*{Test 2 (P1): does the wage gap remain within the same employer?}

\tpart{a}{Hypotheses, written before running the test}
{\small
\begin{tabularx}{\linewidth}{@{}K{3.4cm} Y@{}}
\toprule
$H_0$ & $\beta=0$ \\
$H_1$ & $\beta\neq0$ \\
Significance level & $\alpha=0.05$ family-wise (Holm) \\
Side & two-sided \\
Test & as Test 1, cells = employer $\times$ SOC $\times$ state $\times$ year \\
Smallest difference that matters & 5\% \\
\bottomrule
\end{tabularx}}

Only employers that file both visas for the same job, state and year contribute. If the Test~1
gap vanished here, it would only mean that higher-paying employers hire Australians; if it
remains, the same employer files the two groups differently.

\tpart{b}{Other tests we considered}
Holding the employer fixed removes its pay policy, industry and lawyer, the confounders Test~1
leaves open; fixed effects do this without assuming anything about how employers differ.
{\footnotesize
\begin{tabularx}{\linewidth}{@{}L{2.4cm} Y@{}}
\toprule
\textbf{Also considered} & \textbf{Why we did not use it} \\
\midrule
Paired $t$-test & on cell means it weights a cell with 1 E-3 and 500 H-1B filings like a balanced one. Run as a check with one gap per employer (695, independent): $+5.3\%$, $t=6.59$, $p=9\times10^{-11}$ \\
Wilcoxon signed-rank & the rank version of the paired $t$; tests the median employer gap, not a mean in \%, so it would repeat the paired check \\
Matching E-3 to H-1B filings & discards most H-1B filings for no gain in bias \\
\bottomrule
\end{tabularx}}

\tpart{c}{Result}
{\small
\begin{tabularx}{\linewidth}{@{}K{3.4cm} Y@{}}
\toprule
Rows used & 71{,}060 (3{,}356 E-3), 695 employers, 1{,}906 cells \\
Estimate & $+6.8\%$ ($\hat\beta=0.065$) \\
95\% confidence interval & $+5.3\%$ to $+8.3\%$ \\
Test statistic (d.f.) & $t=9.24$ (694) \\
p-value (raw, Holm-adjusted) & $2.9\times10^{-19}$ ($5.8\times10^{-19}$) \\
Decision & reject $H_0$ \\
Verdict & meaningful: the whole CI is above 5\% \\
\bottomrule
\end{tabularx}}

The same employer files its E-3 workers at a prevailing wage about 7\% higher than its H-1B
workers in the same job, state and year. That is about 75\% of the Test~1 gap, so the premium
is mostly \emph{how} employers file, not \emph{which} employers hire Australians (Test~2 covers
only employers filing both visas, so the two numbers come from different sets of employers).

%------------------------------------------------------------ Test 3
\needspace{18\baselineskip}
\subsection*{Test 3 (P2): are E-3 filings denied more often by employers of the same size?}

\tpart{a}{Hypotheses, written before running the test}
{\small
\begin{tabularx}{\linewidth}{@{}K{3.4cm} Y@{}}
\toprule
$H_0$ & $\beta=0$ \\
$H_1$ & $\beta\neq0$ \\
Significance level & $\alpha=0.05$ family-wise (Holm) \\
Side & two-sided \\
Test & linear probability model, fixed effects for employer-size group $\times$ year, employer-clustered SE \\
Smallest difference that matters & 0.5 percentage points \\
\bottomrule
\end{tabularx}}

$y$ = 1 if denied. Size groups: 1, 2--10, 11--100, 101--1000, $>$1000 filings a year; all six
years are used. 0.5\,pp, because the H-1B denial rate is 0.78\% (exploration set), so it means a $\ge$1.6$\times$
higher risk, and it is over half the whole FY2017--22 fall in the denial rate (1.36\% $\to$ 0.49\%).

\tpart{b}{Other tests we considered}
Employer size is the main confounder (Fig.~\ref{fig:volume}), so the comparison is within size
group and year. A linear probability model gives the gap directly in percentage points, the unit of
the MID, and its clustered SE needs no functional form.
{\footnotesize
\begin{tabularx}{\linewidth}{@{}L{2.4cm} Y@{}}
\toprule
\textbf{Also considered} & \textbf{Why we did not use it} \\
\midrule
Cochran--Mantel--Haenszel test & the classic stratified 2$\times$2 test, but assumes independent filings; the logit gives the same odds ratio with clustering \\
Logistic regression & gives an odds ratio, not pp. Run as a check, same fixed effects and clustering: OR 3.80, $z=24.9$ \\
$\chi^2$ test on the raw 2$\times$2 table & ignores the size confounder the question is about \\
\bottomrule
\end{tabularx}}

\tpart{c}{Result}
{\small
\begin{tabularx}{\linewidth}{@{}K{3.4cm} Y@{}}
\toprule
Rows used & 791{,}227 (14{,}587 E-3), 32{,}956 employers \\
Estimate & $+3.79$\,pp (E-3 5.32\%, H-1B 0.77\% raw; raw gap $+4.55$\,pp) \\
95\% confidence interval & $+3.26$ to $+4.31$\,pp \\
Test statistic (d.f.) & $t=14.12$ (32{,}955) \\
p-value (raw, Holm-adjusted) & $4.0\times10^{-45}$ ($1.6\times10^{-44}$) \\
Decision & reject $H_0$ \\
Verdict & meaningful: the whole CI is above 0.5\,pp \\
\bottomrule
\end{tabularx}}

Among employers of the same size in the same year, E-3 filings are denied 3.8\,pp more often;
employer size explains only 17\% of the raw gap: it shrinks the gap but doesn't reverse it, so
there is no Simpson's paradox. Planned power was low (0.16) because a gap of only 0.5\,pp is
hard to see next to the noise, but the actual gap is about eight times that, so it was detected
easily.
\emph{Exploratory, not planned in advance:}
the gap shrinks from $+6.6$\,pp in FY2018 to $+1.6$\,pp in FY2022 (Fig.~\ref{fig:year}).

%------------------------------------------------------------ Test 4
\needspace{18\baselineskip}
\subsection*{Test 4 (P2): does the denial gap remain within the same employer?}

\tpart{a}{Hypotheses, written before running the test}
{\small
\begin{tabularx}{\linewidth}{@{}K{3.4cm} Y@{}}
\toprule
$H_0$ & $\beta=0$ \\
$H_1$ & $\beta\neq0$ \\
Significance level & $\alpha=0.05$ family-wise (Holm) \\
Side & two-sided \\
Test & as Test 3, cells = employer $\times$ year \\
Smallest difference that matters & 0.5 percentage points \\
\bottomrule
\end{tabularx}}

Only employers filing both visas in the same year contribute. Holding the employer fixed removes
its experience with the process and its form-filling practice, which size groups only approximate.

\tpart{b}{Other tests we considered}
{\footnotesize
\begin{tabularx}{\linewidth}{@{}L{2.4cm} Y@{}}
\toprule
\textbf{Also considered} & \textbf{Why we did not use it} \\
\midrule
CMH test by employer $\times$ year & the standard test for a 2$\times$2 gap within many groups, but gives an odds ratio, not pp. Run as a check: OR 1.89, $p=7.6\times10^{-6}$ \\
Employer-level paired $t$ & one gap per employer (1{,}547), no regression. Run as a check: $+0.59$\,pp, $p=0.018$ \\
Wilcoxon signed-rank & only 306 of 1{,}547 employers have a non-zero gap, and it would repeat the paired $t$ \\
Logistic regression with employer dummies & 1{,}547 dummies but most employer-years have no denial at all, so the odds ratios come out biased \\
\bottomrule
\end{tabularx}}

\tpart{c}{Result}
{\small
\begin{tabularx}{\linewidth}{@{}K{3.4cm} Y@{}}
\toprule
Rows used & 382{,}911 (8{,}003 E-3), 1{,}547 employers \\
Estimate & $+0.51$\,pp (H-1B rate at these employers 0.30\%) \\
95\% confidence interval & $+0.22$ to $+0.81$\,pp \\
Test statistic (d.f.) & $t=3.44$ (1{,}546) \\
p-value (raw, Holm-adjusted) & $5.9\times10^{-4}$ ($5.9\times10^{-4}$) \\
Decision & reject $H_0$ \\
Verdict & possibly too small: the CI includes 0.5\,pp \\
\bottomrule
\end{tabularx}}

Within the same employer the gap is real but small in absolute terms; we can't tell whether it
reaches the 0.5\,pp that matters. Relative to these employers' low H-1B rate (0.30\%) it is about
2.7$\times$. Holding the employer fixed removes 86\% of the Test~3 gap (different sets of employers, as in
Test~2), so most of the E-3 denial
gap comes from \emph{which} employers file E-3, not from how one employer's filings fare.

%------------------------------------------------------------ comparison
\subsection*{Main tests against other families of tests}
All four main tests are the same regression, so we check each with tests from other families.
The \emph{bootstrap} simulates 1{,}999 datasets in which E-3 has no effect, changing whole
employers at a time, and counts how often their $t$ is as large as ours: a $p$-value that doesn't
rely on the $t$ table. We keep one check per family; rank tests were dropped because they
repeat the $t$-based checks. The employer-level checks for Tests 2 and 4 were added after the
planned run.
{\footnotesize
\begin{tabularx}{\linewidth}{@{}L{0.5cm} L{3.3cm} Y L{0.8cm}@{}}
\toprule
\textbf{Test} & \textbf{Method (family)} & \textbf{Result} & \textbf{Agrees} \\
\midrule
T1 & FE regression, clustered $t$ (main) & $+9.1\%$, $p=3\times10^{-42}$ & -- \\
   & bootstrap over employers (resampling) & $p\le0.0005$ & yes \\
   & Welch $t$ on relative wage (two-sample) & $+5.7\%$, $p=2\times10^{-93}$ & yes \\
   & without largest employer & $+9.2\%$, $p=4\times10^{-42}$ & yes \\
\midrule
T2 & FE regression (main) & $+6.8\%$, $p=3\times10^{-19}$ & -- \\
   & bootstrap over employers (resampling) & $p\le0.0005$ & yes \\
   & paired $t$, one gap per employer & $+5.3\%$, $p=9\times10^{-11}$ & yes \\
   & without largest employer & $+6.7\%$, $p=10^{-17}$ & yes \\
\midrule
T3 & FE linear probability (main) & $+3.79$\,pp, $p=4\times10^{-45}$ & -- \\
   & bootstrap over employers (resampling) & $p\le0.0005$ & yes \\
   & logit, same FE, clustered (GLM) & OR 3.80, $p=3\times10^{-137}$ & yes \\
   & without largest employer & $+3.82$\,pp, $p=2\times10^{-45}$ & yes \\
\midrule
T4 & FE linear probability (main) & $+0.51$\,pp, $p=6\times10^{-4}$ & -- \\
   & bootstrap over employers (resampling) & $p\le0.0005$ & yes \\
   & CMH by employer $\times$ year (contingency) & OR 1.89, $p=8\times10^{-6}$ & yes \\
   & paired $t$, one gap per employer & $+0.59$\,pp, $p=0.018$ & yes \\
   & without largest employer & $+0.51$\,pp, $p=9\times10^{-4}$ & yes \\
\bottomrule
\end{tabularx}}

Every alternative gives the same direction and decision. The Welch test's smaller gap ($+5.7\%$)
measures a median-based relative wage, not the mean gap, and the employer-level tests weight each
employer equally, so a typical employer's gap is smaller than the filing-weighted one; Test 4's
paired $t$ ($p=0.018$) is the weakest result, matching its ``possibly too small'' verdict.


%-----------------------------------------------------------------------
\section{Limitations}
\begin{itemize}
  \item \textbf{Association, not cause} (Section~4): we show that filings differ, not that the
  visa causes it.
  \item \textbf{Prevailing wage, not wage paid.} It is the legal minimum, and the wage-level
  column was cut from this extract; relative wage is an indirect proxy (Fig.~\ref{fig:senior}).
  State is coarser than the metro area DOL actually uses.
  \item \textbf{FY2019 left out of the wage tests}, since its file has no prevailing wage, and
  most of its worksites are blank.
  \item \textbf{Duplicates kept.} If some identical rows are export errors, big employers are
  over-weighted; clustering keeps the SEs honest, but we didn't re-run the tests on de-duplicated
  data.
  \item \textbf{Tests 2 and 4 cover only employers filing both visas}, which are larger than average.
  \item \textbf{Denial here is DOL's check of the form}, not the visa decision, which USCIS makes later.
  \item \textbf{Exploratory extras.} The year-by-year results and the employer-level checks of
  Tests 2 and 4 were not planned in advance and are not Holm-corrected. The full test plan, and what changed
  in it, is in \code{report/preregistration.md}.
\end{itemize}

\end{document}
